\chapter{Prepare Your Dev Environment}
\label{chap:fpgalix-devenv}

\section{Introduction}

This chapter covers how to set up the host used for \prjname{FPGAlix}'s hardware and software development. It walks through installing the required software and where it lives on disk, then setting up the environment variables needed to reach it, and finally getting a working \term{USB-Blaster} and serial console access.

Hardware debugging\footnote{Developing inside a \term{VM} is discouraged for this: forwarding the \term{USB-Blaster} through a hypervisor's \term{USB} redirection is a long-standing, widely reported source of \term{JTAG} failures.}, when needed, can only be done through the \term{USB-Blaster}'s \term{JTAG} interface, integrated on the development board itself, using \term{SignalTap}'s in-system logic analyzer from \term{Quartus}. Software debugging, in contrast, can be done either through that same interface or over the network, with \cmdpath{gdbserver} running on the board and attached remotely from the host with \cmdpath{gdb}.

\section{Host requirements}

\prjname{FPGAlix} requires an \term{x86\_64} machine running \term{Ubuntu 22.04 LTS}, with \term{Quartus Prime} (including \term{Platform Designer}) and an \term{ARM GNU} toolchain installed and working.

\section{Required software}

\subsection{apt packages}
\label{subsec:fpgalix-apt-packages}

\begin{shellcode}
sudo apt install -y build-essential bison flex libssl-dev libncurses-dev u-boot-tools cpio unzip rsync bc patch git git-lfs curl wget openssh-server iperf3 picocom pkg-config libopencv-dev libgstreamer1.0-dev libgstreamer-plugins-base1.0-dev libgstrtspserver-1.0-dev gstreamer1.0-plugins-base gstreamer1.0-plugins-good gstreamer1.0-plugins-bad python3 python3-numpy python3-opencv ffmpeg default-jdk
\end{shellcode}

\begin{itemize}
    \item |build-essential|, |bison|, |flex|, |libssl-dev|, |libncurses-dev|, |u-boot-tools|: building the \term{Linux} kernel, the \term{U-Boot} bootloader, and its boot script.
    \item |cpio|, |unzip|, |rsync|, |bc|, |patch|: mandatory host packages for \term{Buildroot} itself, not covered by the group above or already present in the stock \term{Ubuntu} installation, like \cmdpath{tar}, \cmdpath{sed}, \cmdpath{find}, \cmdpath{perl}, \cmdpath{gzip}, etc.
    \item |git|, |git-lfs|: cloning this repository and its submodules. No file in the repository is currently tracked with LFS, but the tooling is installed for forward compatibility.
    \item |curl|, |wget|: fetching installers and, on the board itself, general network testing.
    \item |openssh-server|, |iperf3|: verifying \term{SSH} access and bandwidth to the board once it is on the network.
    \item |picocom|: terminal emulator for the board's serial console.
    \item |pkg-config|, |libopencv-dev|, |libgstreamer1.0-dev|, |libgstreamer-plugins-base1.0-dev|, |libgstrtspserver-1.0-dev|: needed to build \repopath{sw/vision_core} natively on the host (\cmdpath{make intel}), which lets its capture/filter/streaming pipeline be exercised without the board, before cross-compiling it (\cmdpath{make arm}).
    \item |gstreamer1.0-plugins-base|, |gstreamer1.0-plugins-good|, |gstreamer1.0-plugins-bad|: the dev packages above only provide headers for linking the \cmd{vision_core} example; running its intel build additionally needs these plugins loaded at runtime.
    \item |python3|, |python3-numpy|, |python3-opencv|: required by \repopath{sw/camera_driver_2/utils/debug/view.py}, the raw-frame viewer used to debug the \term{V4L2} driver.
    \item |default-jdk|: |sopc2dts|, provided as the \repopath{repos/sopc2dts} submodule, is a plain \term{Java} tool that needs to be compiled before it can be run.
    \item |ffmpeg|: provides \cmdpath{ffplay}, used on the host to view the \term{RTSP} video stream produced by \cmd{vision_core}. \term{VLC} was tried first and dropped due to compatibility problems.
\end{itemize}

\subsection{Install the software, and where it lives}
\label{subsec:fpgalix-install-paths}

\begin{table}[htbp]
\centering
\begin{tabularx}{\textwidth}{@{} l l X @{}}
\toprule
\textbf{Software} & \textbf{Version} & \textbf{Install path} \\
\midrule
\term{Quartus Prime} & \texttt{22.1} & \texttt{\seqsplit{\textasciitilde{}/Programs/Altera/22.1/}} \\
\term{ARM GNU} toolchain & \texttt{12.3.Rel1} & \texttt{\seqsplit{\textasciitilde{}/Programs/Arm/toolchain/arm-gnu-toolchain-12.3.rel1-x86\_64-arm-none-linux-gnueabihf/}} \\
\bottomrule
\end{tabularx}
\caption{Required software and installation paths.}
\label{tab:fpgalix-dev-software}
\end{table}

The user's home folder is sufficient for both; a system-wide installation is not required. See Chapter~\ref{chap:fpgalix-board-bring-up} for where to download each of these.

\subsection{Environment variables}

Add to {\catcode`\~=12 \extpath{~/.bashrc}}, matching the install paths from Section~\ref{subsec:fpgalix-install-paths}:

\begin{shellcodepath}
export QUARTUS_ROOTDIR="$HOME/Programs/Altera/22.1/quartus/"
export QSYS_ROOTDIR="${QUARTUS_ROOTDIR}/sopc_builder/bin"
export PATH="${QUARTUS_ROOTDIR}/bin:$PATH"
export PATH="$HOME/Programs/Arm/toolchain/arm-gnu-toolchain-12.3.rel1-x86_64-arm-none-linux-gnueabihf/bin:$PATH"
export PATH="${QSYS_ROOTDIR}:$PATH"
\end{shellcodepath}

Reload with {\catcode`\~=12 \cmdpath{source ~/.bashrc}} (or open a new shell) to apply the changes just made to the file.

Each line prepends to \term{PATH} rather than appending, so these installs are always found first, ahead of anything already on \term{PATH}, whether or not a same-named tool exists elsewhere on the system.

Not on \term{PATH}, and not meant to be: two other \term{ARM} toolchains exist in this project and are always invoked by their full path from a Makefile, never exported globally.

\begin{itemize}
    \item \term{Buildroot} builds its own internal toolchain (\cmdpath{arm-buildroot-linux-gnueabihf-...}) as part of \repopath{repos/buildroot/output/host/bin/}. It is what \repopath{sw/vision_core}'s \cmdpath{make arm} target uses, and it must stay separate from the toolchain above: mixing the two (e.g. linking an object built with one against a library built with the other) produces \term{glibc} \term{ABI} mismatches.
    \item The host's own native \cmdpath{gcc}/\cmdpath{g++} is used for \repopath{sw/vision_core}'s \cmdpath{make intel} target,\footnote{See Chapter~\ref{chap:fpgalix-vision-core}.} which is a completely separate, \term{x86\_64} build of the same program.
\end{itemize}

\subsection{USB-Blaster and serial console access}

Beyond the packages installed in Section~\ref{subsec:fpgalix-apt-packages}, the board requires access to two \term{USB} devices: the \term{USB-Blaster}, used for \term{JTAG} access (\term{SignalTap}), and the \term{UART-to-USB} serial console. By default, \term{udev} restricts both device nodes to root.

Create a \term{udev} rule granting the invoking user's group access to the \term{USB-Blaster}:

\begin{shellcode}
sudo tee /etc/udev/rules.d/51-usbblaster.rules << 'EOF'
SUBSYSTEM=="usb", ATTR{idVendor}=="09fb", ATTR{idProduct}=="6001", MODE="0666", GROUP="plugdev"
SUBSYSTEM=="usb", ATTR{idVendor}=="09fb", ATTR{idProduct}=="6002", MODE="0666", GROUP="plugdev"
SUBSYSTEM=="usb", ATTR{idVendor}=="09fb", ATTR{idProduct}=="6003", MODE="0666", GROUP="plugdev"
EOF
sudo udevadm control --reload-rules
sudo udevadm trigger
sudo usermod -aG plugdev "$USER"
\end{shellcode}

Without this rule, \term{Quartus} reports ``No JTAG hardware'' even with the \term{USB-Blaster} plugged in. Verify it worked with:

\begin{shellcode}
jtagconfig
\end{shellcode}

which should print something like |1) USB-Blaster [...]|.

The serial console needs no equivalent custom rule: the board's \term{UART-to-USB} connector enumerates as \extpath{/dev/ttyUSB0}, and \term{Ubuntu}'s own default rule already assigns every serial device to the \sysid{dialout} group. The user still needs to be added to it:

\begin{shellcode}
sudo usermod -aG dialout "$USER"
\end{shellcode}

Log out and back in for both group changes to take effect.\footnote{Or run \cmdpath{newgrp plugdev} / \cmdpath{newgrp dialout} on the already-open console to apply the change immediately, without logging out.}

\subsection{Editor}

Both the \term{VHDL} sources under \repopath{hw/} and the \term{C/C++} and \term{Python} sources under \repopath{sw/} are written in \term{VS Code}, with the following extensions:

{\footnotesize
\begin{itemize}
    \item \term{VHDL} linting, snippets, project awareness --- \urlpath{teros-technology.teroshdl} (\term{TerosHDL}).
    \item \term{VHDL} syntax highlighting and language support --- \urlpath{rjyoung.vscode-modern-vhdl-support}.
    \item \term{C/C++} \term{IntelliSense} for \texttt{camera\_driver\_2}, \texttt{hexip\_daemon}, and \texttt{vision\_core} --- \urlpath{ms-vscode.cpptools} (+ \urlpath{ms-vscode.cpptools-extension-pack}).
    \item Makefile awareness (every \repopath{sw/} project builds via \cmdpath{make}) --- \urlpath{ms-vscode.makefile-tools}.
    \item editing/debugging \repopath{sw/camera_driver_2/utils/debug/view.py} --- \urlpath{ms-python.python} (+ \term{Pylance}, \term{debugpy}).
    \item viewing the datasheet/manual PDFs under \repopath{docs/} without leaving the editor --- \urlpath{tomoki1207.pdf}.
\end{itemize}
}

\term{VS Code} is suggested, not required: it is only a text editor here, no build or debug integration is used from it, so any editor works equally well.

\section{Getting the installers}

See Chapter~\ref{chap:fpgalix-board-bring-up} for where to get \term{Quartus Prime}, the \term{ARM GNU} toolchain, and the \term{Ubuntu 22.04 LTS} image.
